\documentclass[11pt, a4paper]{article}\usepackage[utf8]{inputenc}\usepackage{amsmath}\usepackage[margin=1in]{geometry}\begin{document}\title{Catatan Kuliah 7: KKT dan Interior-Point}\maketitle
\section{Kondisi KKT} Kondisi KKT adalah fondasi teoritis untuk optimisasi non-linier terkendala. Ia mengombinasikan fungsi objektif dan kendala ke dalam satu fungsi yang disebut Lagrangian. \section{Interior Point Methods} Metode komputasi modern yang sangat cepat untuk menyelesaikan LP dan NLP cembung. Ia berjalan di *dalam* daerah fisibel, tidak menyusuri tepi seperti Simpleks. \end{document}